% Pursuing Stacks (cotes 134-2 to 134-8): archive cote 134-4, site page 3 (PDF page 4), typescript page 91.
% typescript. Draft compared: ps3.tex:99-179.
% AI-transcribed by Claude Opus 5.5 (claude-opus-5-5, adaptive thinking, effort high), Message Batches API, batch ps1 (msgbatch_01Rv5hCaNTtUt4JHx1yqoHPL), 2026-11-12; image review recorded in records/CHECK_DECISIONS_134_3_4.tsv.
\scanpage{134-4/3}
\mspage{(91)}
which is'nt elementary,  in view of having a standart way for describing
a given homotopy type (defined by an object $C$ in (Cat)) by a``model'' in $M$
depending functorially on $C$,  there does'nt seem to be any general process
for finding one,  even without any demand on exactness properties,  ex\over[typed]{c}{c}\over[typed]{e}{e}pt
of course when $M$ is supposed to be elementary;   in this ca\over[typed]{s}{s}e \xout{xx} $M = A^{\hand{\wedge}}$
we get \xout{of course} the functors
\[ (4) \qquad i^* : (\mathrm{Cat}) \hand{\longrightarrow} \over[typed]{A}{A}^{\hand{\wedge}} \]
associated to test functors $A \hand{\to} (\mathrm{Cat})$,  which can be viewed as a particula\unsure{r}
case of (1),  applied to the case $M = (\mathrm{Cat})$.  Using \xout{this} such functors
(4),  we see that the question of finding a modelizing functor
\[ (*) \qquad \hand{\phi} : (\mathrm{Cat}) \hand{\longrightarrow} M \ , \]
for a more or less general $M$,  is tied up with the question of finding
such a functor from an elementary modelizer $A^{\hand{\wedge}}$ into $M$
\[ (**) \qquad \hand{\psi} : A^{\hand{\wedge}} \hand{\longrightarrow} M \ , \]
M\over[typed]{o}{o}re specifically,  if we got a $\hand{\psi}$ ,  we deduce a $\hand{\phi}$ by composing with
$i^*$ in (4) ,  and conversely,  if we got a $\hand{\phi}$ ,  we deduce a $\hand{\psi}$ by composing
with the canonical functor $i_A$ in (2).  Maybe it's unrealistic to expect
modelizing functors (*) or (**) to exist for rather general M. (\over{w}{W}hich \xout{ones} \inserted[typed]{modelizers}
will turn out to be really ``the interesting ones'' will appear in due \over{\unsure{c}}{c}our-%
se presumably...).  T\over[typed]{h}{h}ere is  one interesting case though when we got such
functors,  namely when
\[ M = (\mathrm{Spaces}) \]
is the category of topological spaces, \inserted{and} taking for $\hand{\psi}$ one of the manifold
\xout{species} \inserted[typed]{avators} of ``geometric realization functor'' ,  associated to a \inserted{suitable} functor
\[ \xout{(5)} \qquad r : A \hand{\longrightarrow} (\mathrm{Spaces}) \]
by taking the canonical extension $r_{\hand{!}}$ to $A^{\hand{\wedge}}$,  commuting with direct limits\unsure{.}
This is presisely the \xout{case} ``highly non typical'' case,  when we get a
pair of adjoint functors  $r_{\hand{!}}$ , $r^*$
% Hand-drawn arrows are marked blue; typed objects remain plain.
\noindent $(5)\qquad$ \diagram{\begin{tikzcd} A^{\hand{\wedge}} \ar[r, shift left,color=blue, "\hand{r_!}"] \& (\mathrm{Spaces}) \ar[l, shift left,color=blue, "\hand{r^*}"] \end{tikzcd}}
which are \uline{both} modelizing.  T\over[typed]{H}{h}e situation  here mimicks very closely the
situation of the pair $(i_A, j_A = i_A^* )$ canonically associated to the elementa-%
r\over[typed]{u}{y} modelizer $A^{\hand{\wedge}}$,  with the ``basic modelizer'' (Cat)  being replaced by
(Spaces),  which therefore can be consid\over[typed]{\illegible{1}}{e}red  as another ``basic modelizer''
of sorts.  In this case  the corresponding functor
\[ \xout{(xxxx)} \ (6) \qquad r_{\hand{!}} i^* : (\mathrm{Cat}) \hand{\longrightarrow} (\over[typed]{\mathrm{Spa}}{\mathrm{Spa}}\over{\mathrm{c}}{\mathrm{c}}\mathrm{es}) \]
mimicks  the functor  $i_A i^*$ of (3) (where on the left hand side $M$ is
taken to be just (Cat),  and on the right (Cat) as the basic modelizer is
replaced by \over[typed]{it}{it}'s \hand{``}next best\hand{''} substitu\over[typed]{e}{t}e (Spaces)).   Here as in (3),  the
modelizing functor we got is neither left nor right exact,  it has a ten-%

\margin{3}
